\begin{solution}{63}
    \begin{subsolution}
        部分子 a、b 的动能分别为
        \begin{equation}
            E _ {\mathrm{a}} = x _ {\mathrm{a}} E, \quad E _ {\mathrm{b}} = x _ {\mathrm{b}} E
            \label{eq:1.s63}
        \end{equation}
        它们远大于其静能，所以部分子 a、b 的质量可忽略，其动量大小分别为
        \begin{equation}
            p _ {\mathrm{a}} = \frac {E _ {\mathrm{a}}}{c}, p _ {\mathrm{b}} = - \frac {E _ {\mathrm{b}}}{c}
            \label{eq:2.s63}
        \end{equation}
        负号表明部分子 a、b 的动量方向相反，在同一条直线上。部分子 a、b 对撞湮灭产生新粒子 S，根据能量守恒定律，S 的能量为
        \begin{equation}
            E _ {\mathrm{s}} = E _ {\mathrm{a}} + E _ {\mathrm{b}}
            \label{eq:3.s63}
        \end{equation}
        根据动量守恒定律，S 的动量大小为
        \begin{equation}
            p _ {\mathrm{s}} = p _ {\mathrm{a}} + p _ {\mathrm{b}}
            \label{eq:4.s63}
        \end{equation}
        由自由粒子的动量和能量关系有
        \begin{equation}
            E _ {\mathrm{s}}^{2} - p _ {\mathrm{s}}^{2} c^{2} = m _ {\mathrm{s}}^{2} c^{4}
            \label{eq:5.s63}
        \end{equation}
        由\eqref{eq:1.s63}\eqref{eq:2.s63}\eqref{eq:3.s63}\eqref{eq:4.s63}\eqref{eq:5.s63}式得
        \begin{equation}
            x _ {\mathrm{a}} x _ {\mathrm{b}} = \frac {m _ {\mathrm{s}}^{2} c^{4}}{4 E^{2}}
            \label{eq:6.s63}
        \end{equation}
        由\eqref{eq:6.s63}式和题给数据得
        \begin{equation}
            x _ {\mathrm{a}} x _ {\mathrm{b}} = 0.0051
            \label{eq:7.s63}
        \end{equation}
    \end{subsolution}
    \begin{subsolution}
        在新粒子 S 静止的参考系中，设 S 衰变成的两光子的频率分别为 $\nu_{1}$ 和 $\nu_{2}$ ，它们的能量和动量大小分别为
        \begin{equation}
            E _ {1} = h \nu_{1}, E _ {2} = h \nu_{2}
            \label{eq:8.s63}
        \end{equation}
        \begin{equation}
            p _ {1} = \frac {h \nu_{1}}{c}, p _ {2} = - \frac {h \nu_{2}}{c}
            \label{eq:9.s63}
        \end{equation}
        负号表明两光子的动量方向相反，在同一条直线上。S衰变前后能量和动量守恒，有
        \begin{equation}
            E _ {1} + E _ {2} = m _ {\mathrm{s}} c^{2}
            \label{eq:10.s63}
        \end{equation}
        \begin{equation}
            p _ {1} + p _ {2} = 0
            \label{eq:11.s63}
        \end{equation}
        由\eqref{eq:8.s63}\eqref{eq:9.s63}\eqref{eq:10.s63}\eqref{eq:11.s63}式得
        \begin{equation}
            \nu_{1} = \nu_{2} = \frac {m _ {s} c^{2}}{2 h}
            \label{eq:12.s63}
        \end{equation}
        由\eqref{eq:12.s63}式和题给数据得
        \begin{equation}
            \nu_{1} = \nu_{2} = \frac {1.0 \times 10^{12} \times 1.60 \times 10^{-19} \mathrm{C} \cdot \mathrm{V}}{2 \times 6.63 \times 10^{-34} \mathrm{J} \cdot \mathrm{s}} \approx 1.2 \times 10^{26} \mathrm{Hz}
            \label{eq:13.s63}
        \end{equation}
    \end{subsolution}
    \begin{subsolution}
        在 S 产生时静止的参照系中，对于 S 衰变到两个粒子 A 的过程，设两个粒子 A 的能量、动量大小分别为 $E_{A1}$ 、 $E_{A2}$ 和 $p_{A1}$ 、 $p_{A2}$ ，由能量守恒定律 $m_{S}c^{2}=E_{A1}+E_{A2}$ 和动量守恒定律 $0=p_{A1}-p_{A2}$ ，以及自由粒子的动量和能量关系
        \begin{equation}
            E _ {\mathrm{Ai}}^{2} = p _ {\mathrm{Ai}}^{2} c^{2} + m _ {\mathrm{A}}^{2} c^{4}, \quad i = 1, 2
            \label{eq:14.s63}
        \end{equation}
        可知，每个粒子 A 的动量大小相等 $p_{A} \equiv p_{A1} = p_{A2}$ ，因而每个粒子 A 的能量为
        \begin{equation}
            E _ {\mathrm{A}} \equiv E _ {\mathrm{A1}} = E _ {\mathrm{A2}} = \frac {1}{2} (m _ {\mathrm{S}} c^{2}) = 5.0 \times 10^{2} \mathrm{GeV} >> m _ {\mathrm{A}} c^{2}
            \label{eq:15.s63}
        \end{equation}
        在粒子 A 衰变到两个同频率光子的过程中，由能量守恒得
        \begin{equation}
            2 h \nu = E _ {\mathrm{A}}
            \label{eq:16.s63}
        \end{equation}
        式中 $\nu$ 是每个光子的频率。设两光子动量之间夹角为 $\theta$ ，由动量守恒定律知，在平行和垂直于粒子 A 运动方向上分别有
        \begin{equation}
            \begin{array}{l l}
                \frac {h \nu}{c} \cos \frac {\theta}{2} + \frac {h \nu}{c} \cos \frac {\theta}{2} & = p _ {\mathrm{A}} \\
                \frac {h \nu}{c} \sin \frac {\theta}{2} - \frac {h \nu}{c} \sin \frac {\theta}{2} & = 0
            \end{array}
            \label{eq:17.s63}
        \end{equation}
        由\eqref{eq:14.s63}\eqref{eq:15.s63}\eqref{eq:16.s63}\eqref{eq:17.s63}式得
        \begin{equation}
            \sin \frac {\theta}{2} = \sqrt {1 - \left(\cos \frac {\theta}{2}\right)^{2}} = \frac {2 m _ {\mathrm{A}}}{m _ {\mathrm{S}}} = 0.0020 << 1
            \label{eq:18.s63}
        \end{equation}
        可见，两光子夹角非常小，利用 $\sin \frac{\theta}{2} \approx \frac{\theta}{2}$ ，由\eqref{eq:18.s63}式得
        \begin{equation}
            \theta \approx \frac {4 m _ {\mathrm{A}}}{m _ {\mathrm{S}}} = \frac {4 \times 1.0}{1.0 \times 10^{3}} = 0.0040
            \label{eq:19.s63}
        \end{equation}
    \end{subsolution}
\end{solution}